\documentclass[10pt,a4paper]{article}

%double si on veut une version prof et une version élève. simple si on veut une seule version.
\newcommand{\buildmode}{double}

\InputIfFileExists{version.tex}{}{}
% Version (définie par le script)
\providecommand{\version}{eleve}

\usepackage{feuillecours}

%%%à modifier à chaque fois

\renewcommand{\niveau}{1G}
\renewcommand{\chapitre}{Suites, généralités}
\renewcommand{\numchap}{1}
\renewcommand{\theme}{Suites}

\begin{document}

\vspace{-0.3cm}

\begin{prerequis}
\begin{itemize}
    \item Calculer l'image d'un nombre par une fonction, notation \(f(x)\).
    \item Résoudre une équation ou une inéquation du premier degré.
    \item Placer des points dans un repère à partir d'un tableau de valeurs.
\end{itemize}
\end{prerequis}

\begin{objectifs}
\begin{itemize}
    \item Connaître les différents modes de génération d'une suite numérique.
    \item Savoir étudier le sens de variation d'une suite.
    \item Savoir représenter graphiquement une suite par un nuage de points.
\end{itemize}
\end{objectifs}

\section{Modes de définition d'une suite}

\begin{defi}
    Une \textbf{suite numérique} est liste de nombre indexée par les entiers naturels. Le \(n\)ème nombre de la liste est noté est notée \(u_n\) plutôt que \(u(n)\) : c'est le \textbf{terme de rang \(n\)} (ou terme d'indice \(n\)) de la suite. On note la suite \((u_n)\), ou \((u_n)_{n\in\N}\).
\end{defi}

\begin{rem}
    On distingue \((u_n)\), qui désigne la suite tout entière (la liste de tous ses termes), et \(u_n\), qui désigne un seul \textbf{terme} de cette suite (un nombre). \(u_0\) est le premier terme lorsque la suite est définie à partir de \(n=0\) ; \(u_{n+1}\) désigne le terme qui suit \(u_n\), et \(u_{n-1}\) celui qui le précède.
\end{rem}

\subsection{Définition explicite}

\begin{defi}
    Une suite \((u_n)\) est définie de façon \textbf{explicite} lorsqu'il existe une fonction \(f\) telle que, pour tout entier naturel \(n\), \(u_n = f(n)\).
\end{defi}

\begin{exple}
    La suite \((u_n)\) définie pour tout entier naturel \(n\) par \(u_n = 2n+5\) est définie explicitement, avec \(f(n)=2n+5\).
    \begin{itemize}
        \item \(u_0 = 2\times 0+5 = 5\).
        \item \(u_5 = 2\times 5+5 = 15\).
        \item Pour calculer \(u_{100}\), il suffit de remplacer \(n\) par \(100\) : \(u_{100} = 2\times 100+5 = 205\), sans avoir besoin de connaître \(u_{99}\).
    \end{itemize}
\end{exple}

\begin{met}
    Pour exprimer \(u_{n+1}\) en fonction de \(n\) à partir d'une définition explicite \(u_n=f(n)\), on remplace \(n\) par \(n+1\) \textbf{partout} où il apparaît dans l'expression de \(f(n)\).
\end{met}

\begin{exple}
    Avec \(u_n = 2n+5\), on a \(u_{n+1} = 2(n+1)+5 = 2n+7\).
\end{exple}

\begin{rem}
    Attention à ne pas confondre \(u_{n+1}\) (le terme qui suit \(u_n\)) et \(u_n+1\) (le nombre \(u_n\) auquel on ajoute \(1\)) : ce sont en général deux nombres différents.
\end{rem}

\subsection{Définition par récurrence}

\begin{defi}
    Une suite \((u_n)\) est définie par \textbf{récurrence} lorsqu'on donne :
    \begin{itemize}
        \item son premier terme (par exemple \(u_0\)) ;
        \item une relation permettant de calculer chaque terme à partir du précédent, de la forme \(u_{n+1} = f(u_n)\).
    \end{itemize}
    On la note souvent \(\left\{\begin{array}{rcl} u_0 &=& \ldots \\ u_{n+1} &=& \ldots \end{array}\right.\)
\end{defi}

\begin{exple}
    La suite \((u_n)\) définie par \(u_0=3\) et, pour tout entier naturel \(n\), \(u_{n+1}=2u_n+1\), est définie par récurrence.
    \begin{itemize}
        \item \(u_1 = 2u_0+1 = 2\times 3+1 = 7\).
        \item \(u_2 = 2u_1+1 = 2\times 7+1 = 15\).
    \end{itemize}
\end{exple}

\begin{met}
    Pour une suite définie par récurrence, on ne peut pas calculer directement un terme sans connaître le précédent : pour obtenir \(u_{10}\), il faut calculer successivement \(u_1\), \(u_2\), \ldots, jusqu'à \(u_{10}\), \textbf{de proche en proche}.
\end{met}

\prof{
    Bien insister sur cette différence avec la définition explicite : c'est souvent la principale source de confusion en début de chapitre (\og pourquoi je ne peux pas juste remplacer \(n\) par \(10\) ? \fg{}). On peut aussi signaler, sans le démontrer ici, que deviner une formule explicite à partir de quelques termes seulement ne suffit jamais à la prouver.
}

\section{Autres modes de génération d'une suite}

\begin{rem}
    Une suite numérique peut aussi être définie par d'autres procédés, qui se ramènent en général à une définition explicite ou par récurrence une fois analysés.
\end{rem}

\subsection{À partir d'un motif géométrique}

\begin{met}
    Lorsqu'une suite est définie par une construction géométrique qui se répète (un motif), on procède ainsi :
    \begin{enumerate}
        \item calculer les premiers termes en comptant directement sur la figure ;
        \item observer comment on passe d'une figure à la suivante, pour en déduire une relation de récurrence ;
        \item si besoin, chercher une expression explicite cohérente avec les premiers termes calculés.
    \end{enumerate}
\end{met}

\subsection{À partir d'un algorithme}

\begin{rem}
    Une boucle qui met à jour une variable \(u\) à chaque tour génère une suite \((u_n)\) définie par récurrence : la valeur de \(u\) avant la boucle est \(u_0\), et l'instruction exécutée à chaque tour donne la relation entre \(u_{n+1}\) et \(u_n\).
\end{rem}

\begin{met}
    Pour retrouver la suite définie par un algorithme, on identifie :
    \begin{itemize}
        \item la valeur initiale de la variable (avant la boucle) : c'est \(u_0\) ;
        \item l'instruction exécutée à chaque tour de boucle : elle donne \(u_{n+1}\) en fonction de \(u_n\).
    \end{itemize}
\end{met}

\section{Sens de variation d'une suite}

\begin{defi}
    Soit \((u_n)\) une suite numérique. On dit que \((u_n)\) est :
    \begin{itemize}
        \item \textbf{croissante} si, pour tout entier naturel \(n\), \(u_{n+1}\geq u_n\) ;
        \item \textbf{décroissante} si, pour tout entier naturel \(n\), \(u_{n+1}\leq u_n\) ;
        \item \textbf{constante} si, pour tout entier naturel \(n\), \(u_{n+1} = u_n\).
    \end{itemize}
    Une suite qui n'est ni croissante ni décroissante est dite \textbf{non monotone}.
\end{defi}

\begin{rem}
    Ces définitions doivent être vraies pour \textbf{tout} entier naturel \(n\), et pas seulement pour les premiers termes : vérifier que \(u_0\leq u_1\leq u_2\leq u_3\), par exemple, ne suffit pas à conclure que la suite est croissante.
\end{rem}

\begin{met}
    Pour étudier le sens de variation d'une suite \((u_n)\), on étudie le signe de \(u_{n+1}-u_n\) pour tout entier naturel \(n\) :
    \begin{itemize}
        \item si \(u_{n+1}-u_n\geq 0\) pour tout \(n\), alors \((u_n)\) est croissante ;
        \item si \(u_{n+1}-u_n\leq 0\) pour tout \(n\), alors \((u_n)\) est décroissante.
    \end{itemize}
\end{met}

\begin{exple}
    Soit \((u_n)\) définie pour tout entier naturel \(n\) par \(u_n = 5n-3\). Pour tout entier naturel \(n\) :
    \[
    u_{n+1}-u_n = \big(5(n+1)-3\big)-\big(5n-3\big) = 5.
    \]
    Comme \(5>0\), on a \(u_{n+1}-u_n>0\) pour tout \(n\) : la suite \((u_n)\) est croissante.
\end{exple}

\begin{rem}
    Lorsqu'une suite est définie par récurrence sous la forme \(u_{n+1}=u_n+f(n)\), le signe de \(f(n)\) donne directement le signe de \(u_{n+1}-u_n\) : inutile de recalculer \(u_{n+1}-u_n\) séparément.
\end{rem}

\begin{exple}
    Soit \((u_n)\) définie par \(u_0=1\) et, pour tout entier naturel \(n\), \(u_{n+1}=u_n+2n+3\). Alors \(u_{n+1}-u_n = 2n+3\). Comme \(2n+3>0\) pour tout entier naturel \(n\), la suite \((u_n)\) est croissante.
\end{exple}

\section{Représentation graphique : nuage de points}

\begin{defi}
    Représenter graphiquement une suite \((u_n)\), c'est placer dans un repère les points de coordonnées \((n\ ;\ u_n)\), pour les entiers naturels \(n\) considérés. L'ensemble de ces points est appelé un \textbf{nuage de points}.
\end{defi}

\begin{rem}
    Contrairement à la courbe représentative d'une fonction, les points d'un nuage de points ne sont \textbf{pas reliés entre eux} : une suite n'est définie que pour des valeurs entières de \(n\), et n'a pas de sens \og entre deux entiers \fg{}.
\end{rem}

\begin{met}
    Pour représenter une suite \((u_n)\) :
    \begin{enumerate}
        \item calculer (ou lire) les valeurs de \(u_n\) pour les premiers entiers naturels \(n\), éventuellement dans un tableau de valeurs ;
        \item placer, dans un repère, chaque point de coordonnées \((n\ ;\ u_n)\), sans les relier.
    \end{enumerate}
\end{met}

\begin{exple}
    Voici le nuage de points de la suite \((u_n)\) définie pour tout entier naturel \(n\) par \(u_n = -n^2+6n\), pour \(n\) allant de \(0\) à \(6\).

    \begin{center}
    \begin{tabular}{|c|c|c|c|c|c|c|c|}
        \hline
        \(n\) & 0 & 1 & 2 & 3 & 4 & 5 & 6 \\
        \hline
        \(u_n\) & 0 & 5 & 8 & 9 & 8 & 5 & 0 \\
        \hline
    \end{tabular}
    \end{center}

    \begin{center}
    \begin{tikzpicture}
    \begin{axis}[
        axis lines=middle,
        xmin=-0.5, xmax=6.5,
        ymin=-1, ymax=10,
        xtick={0,1,2,3,4,5,6},
        ytick={0,2,4,6,8},
        xlabel={\(n\)},
        ylabel={\(u_n\)},
        width=9cm,
        height=6cm
    ]
    \addplot[only marks, mark=*, mark size=2pt, blue] coordinates {
        (0,0) (1,5) (2,8) (3,9) (4,8) (5,5) (6,0)
    };
    \end{axis}
    \end{tikzpicture}
    \end{center}
\end{exple}

\begin{rem}
    Un nuage de points permet de \textbf{conjecturer} le sens de variation d'une suite (les points semblent monter ou descendre), mais cette conjecture doit ensuite être démontrée algébriquement, par exemple à l'aide du signe de \(u_{n+1}-u_n\) (partie précédente). Dans l'exemple ci-dessus, la suite semble croissante jusqu'à \(n=3\), puis décroissante : elle n'est pas monotone.
\end{rem}

\prof{
    On peut ouvrir, sans le développer, sur le fait que les termes d'une suite peuvent sembler se rapprocher d'une valeur sans jamais l'atteindre (idée intuitive de limite, hors programme à ce stade).
}

\end{document}
